\section{Implementation, interfaces, and trust boundaries}
\label{sec:implementation}

The implementation adds three production modules and changes the existing
orbit producer at its periodic-closure call site.  The existing independent
orbit verifier is retained unchanged.  There is no new dependency in the
production path: the modules use the Python standard library and the
repository's existing geometry and integer codecs.  Regina is used only in
an independent experimental audit.

\subsection{A compatible scheduling interface}

The orbit producer accepts
\begin{verbatim}
count_orbits(n, pairings, ..., merger_scheduler='adaptive')
\end{verbatim}
with three scheduling choices.  \code{legacy} performs the literal historical
restart scan.  \code{queue} uses the event queue from the start of every
periodic closure.  The default, \code{adaptive}, spends a bounded scan
allowance before switching to that queue.  Both existing periodic rules,
\code{fine\_wilf} and \code{aht}, remain available.

The scheduling module receives the merger as a callback; it contains no
second implementation of periodic mathematics.  The callback must be
deterministic and side-effect free.  It returns the same replacement pairing
that the old scan would have installed.  Candidate ordering and generation
invalidation then account for all differences in execution.

The certificate format does not change.  Queue entries use stable original
row identifiers internally, but a successful operation records the current
list ranks expected by the existing verifier.  The optional statistics add
queue runs, adaptive switches, peak queue occupancy, and stale pops.  These
are execution diagnostics, not parts of the mathematical certificate.

The adaptive implementation has a small fast path: it obtains the first two
periodic rows lazily.  If their merger succeeds, it restarts immediately.
Only after that first failure does it finish collecting the periodic index
list for the remaining lexicographic scan.  Consequently an easy sequence
of first-pair successes does not rebuild a full index list at each merge.
The first test is charged to the same allowance and is not repeated.  The
exact simulation theorem and the quadratic test bound are unchanged.

Callbacks are checked throughout candidate discovery, queue updates, stale
pops, rank conversion, and output construction.  A raised exception
propagates.  Python allocation, integer arithmetic, and heap primitives are
not hard-real-time preemptible; cooperative checking must not be described
as a strict elapsed-time bound.

\subsection{Weighted profiles and reusable traces}

The new module \path{fastunknot/weighted_orbits.py} exports
\begin{verbatim}
weighted_orbit_profile(n, pairings, runs, ...)
replay_weighted_orbit_profile(n, pairings, runs, certificate, ...)
\end{verbatim}
where \code{runs} is the full ordered half-open partition into triples
\code{(start, stop, weight\_tuple)}.  All tuples have a common positive
dimension; a separate \code{dimension} argument resolves the empty-input
case.  The result contains completion status, orbit count, sorted profile
groups, dimension, producer cycle count when available, and diagnostics.
Each group is \code{(multiplicity, weight\_tuple)}.  Requesting
\code{record\_certificate=True} returns the checked unweighted trace.

Supplying an existing certificate avoids producer scheduling.  The replay
function first validates that trace through
\path{fastunknot/interval_orbit_verify.py}, then performs the weight
transfers.  It does not trust a claimed terminal count or a claimed set of
weights embedded in an arbitrary trace.  A trace is intentionally reusable
for different input weight fields on the same ordered interval system.

Invalid weight partitions and invalid traces raise
\code{WeightedOrbitError}.  Caller cancellation exceptions propagate even
when their type is a common validation exception such as \code{ValueError}.
A cycle-limited producer that cannot finish returns neither an orbit count
nor a partial profile nor a certificate.  These semantics prevent an
interrupted computation from becoming negative evidence.

\subsection{A supplied-surface component query}

The module \path{fastunknot/normal_component_profile.py} exports
\begin{verbatim}
normal_component_profile(triangulation, coordinates,
                         record_certificate=True)
verify_normal_component_profile(triangulation, coordinates,
                                certificate)
\end{verbatim}
as well as the component-weight construction for research use.  A completed
query reports the following summary fields:
\begin{center}
\begin{tabular}{lp{0.43\textwidth}}
\toprule
\code{components} & Total connected components. \\
\code{disk\_components} & All disk components. \\
\code{compressing\_disk\_components} & Disks with essential boundary. \\
\code{has\_compressing\_disk\_component} & Whether that count is positive. \\
\bottomrule
\end{tabular}
\end{center}
Its group records also report the five integer fields, the two boundary
parities, and both disk predicates.

The previous topology interface retains its whole-surface
\code{compressing\_disk} field.  That field and the new existential
component field answer different questions.  Code consuming either result
should name that quantifier explicitly.  For example, a meridian plus a
vertex-linking disk is not a connected compressing disk, but it does contain
a compressing-disk component.  The coordinate vector can have greatest
common divisor one, so dividing a common coordinate factor would not solve
this problem.

The component certificate contains an algorithm tag, a source fingerprint,
the unweighted orbit certificate, and the profile multiset.  The verifier
reconstructs the geometry and weights from the supplied triangulation and
normal vector, validates the trace independently of its producer, replays
weights, and compares the claimed profiles.  Large signed integers survive
the existing hexadecimal JSON transport.  A fingerprint is only an identity
check: it cannot replace the matching equations, cellular construction, or
local mathematical verification.

This verifier shares the geometry and initial-weight construction with its
producer.  The independent Regina and literal-component comparisons in
\cref{sec:experiments} test that shared boundary.  They do not make the
checker a separately developed geometric implementation.  A future formal
verification effort should preserve this distinction between independent
trace checking and shared geometric preprocessing.

\subsection{What is deliberately outside the returned claim}

The query validates the topology of its input manifold, but it does not
certify that this manifold is the exterior of a particular planar diagram.
An essential disk in a verified knot exterior would certify the unknot;
the missing link is the source-preserving construction of that exterior.
The ordinary recognition dispatcher therefore receives no new unconditional
acceptance branch in this patch.

A negative query means that the supplied normal surface has no compressing
disk among its components.  It is not a certificate that all normal vectors
fail, that the boundary is incompressible, or that a knot is nontrivial.
Likewise, a positive profile counts disk components but does not yet give
the chosen disk's complete normal coordinates, attaching maps for a cut,
or an explicit isotopy.  These stronger artifacts require additional work
identified in \cref{sec:research}.

The implementation retains a complete orbit trace before weighted replay.
Its representation is polynomial under the established AHT trace bound,
but storage can be material.  Streaming verified operations is a useful
future memory improvement; it must preserve the ability to check a saved
certificate independently.  The event queue can also have quadratically
many entries.  Neither cost is hidden behind the word ``compressed.''
